% Part: many-valued-logic
% Chapter: three-valued-logics
% Section: kleene

\documentclass[../../../include/open-logic-section]{subfiles}

\begin{document}

\olfileid{mvl}{thr}{skl}

\olsection{క్లీని తర్కాలు}

గణన ప్రక్రియల ఫలితాలుగా సత్యమూల్యాలను భావించే
దృష్టితో స్టీఫెన్ క్లీని రెండు మూడు-విలువల తర్కాలను
ప్రవేశపెట్టాడు. ఒక ప్రక్రియ $\True$ లేదా $\False$
ఫలితాన్ని ఇవ్వవచ్చు; కానీ అది ముగియకపోవచ్చూ.
అప్పుడు సంబంధిత సత్యమూల్యం నిర్వచితం కాదు;
దానిని $\Undef$ అనే సత్యమూల్యంతో సూచిస్తాం.

ఒక ప్రతిపాదన~$!A$ నిషేధాన్ని గణించాలంటే, మొదట
$!A$ విలువను గణించి, తరువాత దానికి వ్యతిరేకమైన
ఫలితాన్ని ఇవ్వాలి. $!A$ గణన ముగియకపోతే మొత్తం
ప్రక్రియ కూడా ముగియదు; కాబట్టి $\Undef$ నిషేధం
$\Undef$.

సంయోగం $!A \land !B$ను గణించడానికి రెండు
మార్గాలున్నాయి. మొదట~$!A$ను, తరువాత $!B$ను
గణించి, రెండింటి ఫలితాలూ~$\True$ అయితే
$\True$ను, లేకపోతే $\False$ను ఇవ్వవచ్చు.
ఏ గణన అయినా ఆగకపోతే మొత్తం ప్రక్రియ కూడా
ఆగదు. ఈ పద్ధతిలో ఒక సంయోగ భాగం నిర్వచితం
కాకపోతే సంయోగమూ నిర్వచితం కాదు. వికల్పానికీ
ఇదే వర్తిస్తుంది.

అయితే $!A$, $!B$లను సమాంతరంగా గణించగలిగితే
మరింత చెప్పవచ్చు. రెండు ప్రక్రియల్లో ఒకటి ఆగి
$\False$ ఇస్తే, సంయోగం తప్పకుండా అబద్ధమని
తెలిసి అక్కడే ఆగవచ్చు. అందువల్ల ఒక భాగం
నిర్వచితం కాకపోయినా, మరొక భాగం అబద్ధమైతే
సంయోగం అబద్ధమే. అలాగే వికల్పాన్ని సమాంతరంగా
గణించేటప్పుడు ఒక భాగం నిజమని తేలగానే,
వికల్పం నిజమని చెప్పి ఆగవచ్చు. ఈ సందర్భంలో,
ఒక భాగం నిర్వచితం కాకపోయినా సంయుక్త వాక్యపు
ఫలితం తెలియవచ్చు. ఈ వ్యాఖ్యానంలో $\Undef$ను
“నిర్వచితం కానిది” కన్నా “తెలియనిది”గా
చదవవచ్చు.

ఈ రెండు వ్యాఖ్యానాల నుంచి క్లీని యొక్క బలమైన,
బలహీనమైన తర్కాలు వస్తాయి. సోపాధిక సంయోజకాన్ని
$\lnot !A \lor !B$కు సమానార్థకంగా నిర్వచిస్తాం.

\begin{defn}
\emph{బలమైన క్లీని తర్కం}~$\LogKs$ను కింది మాత్రికతో
నిర్వచిస్తాం:
\begin{enumerate}
  \item చరాల నుంచి $\lnot$, $\land$, $\lor$, $\lif$
  సంయోజకాలతో మాత్రమే నిర్మించిన ప్రామాణిక
  ప్రతిజ్ఞావాక్య భాష $\Lang L_0$ యొక్క ఉపభాష.
  \item సత్యమూల్యాల సమితి $V = \{\True, \Undef, \False\}$.
  \item $\True$ మాత్రమే నిర్దేశిత విలువ; అంటే
  $V^+ = \{\True\}$.
  \item సత్యమూల్య ప్రమేయాలను కింది పట్టికలు ఇస్తాయి:
  \begin{center}
    \begin{tabular}{c|c}
      $\tf{\lnot}$ & \\
      \hline
      $\True$ & $\False$ \\
      $\Undef$ & $\Undef$ \\
      $\False$ & $\True$
    \end{tabular}
    \quad
    \begin{tabular}{c|ccc}
      $\tf{\land}[\LogKs]$ & $\True$ & $\Undef$ & $\False$ \\
      \hline
      $\True$ & $\True$ & $\Undef$ & $\False$ \\
      $\Undef$ & $\Undef$ & $\Undef$ & $\False$\\
      $\False$ & $\False$ & $\False$ & $\False$
    \end{tabular}
    \\[2ex]
    \begin{tabular}{c|ccc}
      $\tf{\lor}[\LogKs]$ & $\True$ & $\Undef$ & $\False$ \\
      \hline
      $\True$ & $\True$ & $\True$ & $\True$ \\
      $\Undef$ & $\True$ & $\Undef$ & $\Undef$ \\
      $\False$ & $\True$ & $\Undef$ & $\False$
    \end{tabular}
    \quad
    \begin{tabular}{c|ccc}
      $\tf{\lif}[\LogKs]$ & $\True$ & $\Undef$ & $\False$ \\
      \hline
      $\True$ & $\True$ & $\Undef$ & $\False$ \\
      $\Undef$ & $\True$ & $\Undef$ & $\Undef$  \\
      $\False$ & $\True$ & $\True$ & $\True$
    \end{tabular}
  \end{center}
\end{enumerate}
\sourcecorrection{OLTEMVLKLE-001}{ఆంగ్ల మూలం
ప్రామాణిక భాషను పేర్కొన్నా, అందులోని అసత్య
స్థిరాంకానికి విలువ ఇవ్వలేదు. తరువాతి
సర్వసత్యాలు లేవనే నిరూపణ నిలవాలంటే,
ఈ మాత్రికను చరాలు మరియు పట్టికల్లోని నాలుగు
సంయోజకాలతో ఏర్పడే ఉపభాషకే పరిమితం చేయాలి;
అసత్య స్థిరాంకం ఇందులో లేదు.}
\end{defn}

\begin{defn}
\emph{బలహీనమైన క్లీని తర్కం}~$\LogKw$ను కింది మాత్రికతో
నిర్వచిస్తాం:
\begin{enumerate}
  \item చరాల నుంచి $\lnot$, $\land$, $\lor$, $\lif$
  సంయోజకాలతో మాత్రమే నిర్మించిన ప్రామాణిక
  ప్రతిజ్ఞావాక్య భాష $\Lang L_0$ యొక్క ఉపభాష.
  \item సత్యమూల్యాల సమితి $V = \{\True, \Undef, \False\}$.
  \item $\True$ మాత్రమే నిర్దేశిత విలువ; అంటే
  $V^+ = \{\True\}$.
  \item సత్యమూల్య ప్రమేయాలను కింది పట్టికలు ఇస్తాయి:
  \begin{center}
    \begin{tabular}{c|c}
      $\tf{\lnot}$ & \\
      \hline
      $\True$ & $\False$ \\
      $\Undef$ & $\Undef$ \\
      $\False$ & $\True$
    \end{tabular}
    \quad
    \begin{tabular}{c|ccc}
      $\tf{\land}[\LogKw]$ & $\True$ & $\Undef$ & $\False$ \\
      \hline
      $\True$ & $\True$ & $\Undef$ & $\False$ \\
      $\Undef$ & $\Undef$ & $\Undef$ & $\Undef$\\
      $\False$ & $\False$ & $\Undef$ & $\False$
    \end{tabular}
    \\[2ex]
    \begin{tabular}{c|ccc}
      $\tf{\lor}[\LogKw]$ & $\True$ & $\Undef$ & $\False$ \\
      \hline
      $\True$ & $\True$ & $\Undef$ & $\True$ \\
      $\Undef$ & $\Undef$ & $\Undef$ & $\Undef$ \\
      $\False$ & $\True$ & $\Undef$ & $\False$
    \end{tabular}
    \quad
    \begin{tabular}{c|ccc}
      $\tf{\lif}[\LogKw]$ & $\True$ & $\Undef$ & $\False$ \\
      \hline
      $\True$ & $\True$ & $\Undef$ & $\False$ \\
      $\Undef$ & $\Undef$ & $\Undef$ & $\Undef$  \\
      $\False$ & $\True$ & $\Undef$ & $\True$
    \end{tabular}
  \end{center}
\end{enumerate}
\sourcecorrection{OLTEMVLKLE-002}{ఇక్కడ కూడా ఆంగ్ల
మూలం ప్రామాణిక భాషను పేర్కొని అసత్య స్థిరాంకపు
విలువను ఇవ్వలేదు. కింది నిరూపణకు అనుగుణంగా
చరాలు, పట్టికల్లోని నాలుగు సంయోజకాలతో ఏర్పడే
ఉపభాషనే తీసుకున్నాం; అసత్య స్థిరాంకాన్ని
చేర్చలేదు.}
\end{defn}

\begin{prop}
  $\LogKs$, $\LogKw$లలో సర్వసత్యాలు లేవు.
\end{prop}

\begin{proof}
  అన్ని \tetoken{ప్రతిజ్ఞావాక్య చరాలకు}{!!{propositional variable}s}~$p$
  $\pAssign v(p) = \Undef$ అయితే, ఏ
  సూత్రం~$!A$ సత్యమూల్యమైనా~$\pValue v(!A) =
  \Undef$ అవుతుంది. ఎందుకంటే రెండు తర్కాల్లోనూ
  \[
    \tf{\lnot}(\Undef) = \tf{\lor}(\Undef, \Undef) = \tf{\land}(\Undef,
  \Undef) = \tf{\lif}(\Undef, \Undef) = \Undef
  \]
  అవుతుంది. $\LogKs$, $\LogKw$లలోనూ
  $\Undef \notin V^+$ కనుక, ఈ
  \tetoken{సత్యమూల్య కేటాయింపులో}{!!{valuation}}
  $!A$ నిర్దేశిత విలువను పొందదు.
\end{proof}

బలహీన, బలమైన క్లీని తర్కాలు రెండింటిలోనూ
సర్వసత్యాలు లేకపోయినా, వాటి పర్యవసాన
సంబంధాలు అల్పమైనవి కావు.

\begin{prob}
  కింది సంబంధాల్లో ఏవి (a) బలమైన,
  (b) బలహీనమైన క్లీని తర్కంలో నిలుస్తాయి?
  ప్రతిదానికీ సత్య పట్టిక ఇవ్వండి.
  \begin{enumerate}
    \item $p, p \lif q \Entails q$
    \item $p \lor q, \lnot p \Entails q$
    \item $p \land q \Entails p$
    \item $p \Entails p \land p$
    \item $p \Entails p \lor q$
  \end{enumerate}
\end{prob}

ద్మిత్రి బోచ్వార్ $\Undef$ను “అర్థరహితం”గా
వ్యాఖ్యానించాడు. వైరుధ్యపూరిత వాక్యాలకు
$\Undef$ విలువను నిర్దేశించడం ద్వారా
అసత్యవాది వైరుధ్యాభాసం వంటి వాటిని
పరిష్కరించడానికి ప్రయత్నించాడు. అతడు,
ముఖ్యంగా బలహీన క్లీని తర్కాన్ని అదనపు
సంయోజకాలతో విస్తరించిన తర్కాన్ని ప్రవేశపెట్టాడు.
వాటిలో రెండు “బాహ్య నిషేధం”, “నిర్వచితం
కానిది” అనే సంచాలకాలు:
\begin{center}
  \begin{tabular}{c|c}
    $\tf{\sim}$ & \\
    \hline
    $\True$ & $\False$ \\
    $\Undef$ & $\True$ \\
    $\False$ & $\True$
  \end{tabular}
\quad
  \begin{tabular}{c|c}
    $\tf{+}$ & \\
    \hline
    $\True$ & $\False$ \\
    $\Undef$ & $\True$ \\
    $\False$ & $\False$
  \end{tabular}
\end{center}

\begin{prob}
  బోచ్వార్ తర్కంలో $\sim$ను $\lnot$,
  $+$లతో నిర్వచించగలరా? అంటే $\sim$ లేకుండా,
  \tetoken{ప్రతిజ్ఞావాక్య చరం}{!!{propositional
  variable}}~$p$ మాత్రమే కలిగి, ఎప్పుడూ
  $\mathord{\sim}p$తో సమాన సత్యమూల్యాన్ని
  పొందే \tetoken{సూత్రాన్ని}{!!a{formula}}
  కనుగొనండి. మీ సమాధానాన్ని సమర్థించే
  సత్య పట్టిక ఇవ్వండి.
\end{prob}

\end{document}
